\section{基于KG-GCN的燃煤机组汽轮机系统故障诊断方法}
\label{sec:kg_gcn_v3}

\subsection{引言}

第二章基于正常运行DCS数据建立了设备级GRU状态监测模型，使模型能够在给定运行边界和操作状态下给出汽轮机设备的健康状态参考；第三章进一步利用设备IO、DCS点表以及公开故障资料，将汽轮机主通流、再热、抽汽回热、给水及冷端系统中的设备、参数、异常和故障关系组织为领域知识图谱。前者描述设备在正常工况下“应该处于什么状态”，后者描述不同设备和状态参数之间“为什么存在联系”。对于复杂变工况下的汽轮机故障诊断，仅使用原始DCS参数容易将正常负荷变化与故障偏差混合，而仅依靠静态规则又难以表达当前时刻多参数的动态异常程度，因此有必要将健康状态偏差与知识拓扑进一步融合。

图卷积网络（Graph Convolutional Network, GCN）能够在非欧几里得图结构上实现邻域信息聚合，使一个参数节点的表示同时包含自身状态及其结构邻域状态\cite{Kipf2017GCN}。这一机制与汽轮机故障沿主通流、抽汽和再热接口传播的物理过程具有较好的对应关系。本文在课题组KG-GCN技术路线\cite{YangHaodong2026CFB}基础上，将第三章经过证据审计的汽轮机知识关系投影为参数图，将第二章GRU健康状态模型产生的实测--预测偏差映射到图节点，并采用GCN提取多参数关联特征。

与锅炉主辅机跨多个设备系统的诊断任务不同，本文第五章选择VHP作为代表性算例，其24个DCS状态点固定对应VHP本体、一级抽汽支路和一次再热接口等明确物理位置。初步结构实验表明，若直接对全部节点采用全局平均池化，节点身份信息会被明显削弱，不利于区分“异常发生于VHP本体、一级抽汽支路还是再热接口”等局部传播模式。因此，本文在保持GCN邻域聚合机制不变的基础上，采用保留节点顺序的node-aware图级读出方式形成分类特征。该修改属于面向固定测点参数子图的工程适配，模型仍以知识图谱邻接矩阵约束图卷积传播。

\subsection{图卷积网络与参数图构建}
\label{subsec:gcn_v3}

设当前诊断对象对应的参数图为
\begin{equation}
    G_p=(V_p,E_p),
\end{equation}
其中$V_p=\{v_1,v_2,\ldots,v_N\}$表示能够加载实时数值的参数节点集合，$E_p$表示由汽轮机系统知识图谱投影得到的参数关联边。第三章完整知识图谱同时包含设备、故障、异常及处置等语义实体，但GCN实时输入节点必须能够与DCS状态量建立明确映射，因此需要从系统级知识图谱中抽取诊断对象对应的参数子图。

考虑到本文缺乏现场真实故障样本，第五章采用半物理故障场景进行方法验证。为避免“利用同一故障知识既生成样本又定义分类拓扑”造成循环论证，主实验中的KG邻接矩阵仅采用A0级本机设备拓扑、参数归属和直接热力接口关系进行投影；A1级故障机理关系主要用于第三章知识解释及第五章故障场景机理审计，不直接作为主分类邻接矩阵中的故障答案边。由此，模型获得的是系统结构先验，而不是故障类别标签的隐式编码。

若节点$v_i$与$v_j$之间存在允许用于图卷积传播的知识连接，则邻接矩阵$A\in\mathbb R^{N\times N}$中相应元素取1。为保留节点自身特征，引入单位矩阵$I_N$构造增强邻接矩阵
\begin{equation}
    \tilde A=A+I_N.
\end{equation}
定义度矩阵
\begin{equation}
    \tilde D_{ii}=\sum_j\tilde A_{ij},
\end{equation}
得到对称归一化邻接矩阵
\begin{equation}
    \bar A=\tilde D^{-\frac{1}{2}}\tilde A\tilde D^{-\frac{1}{2}}.
\end{equation}
设第$l$层节点特征矩阵为$H^{(l)}\in\mathbb R^{N\times d_l}$，则GCN层间传播可写为
\begin{equation}
    H^{(l+1)}=\sigma\left(\bar A H^{(l)}W^{(l)}\right),
\end{equation}
其中$W^{(l)}$为可学习权重矩阵，$\sigma(\cdot)$为非线性激活函数。该过程使参数节点仅从具有真实设备拓扑或热力接口关系的邻域中接收特征，从而将系统结构知识嵌入多参数动态特征提取过程。

\subsection{KG-GCN模型输入构建}
\label{subsec:kg_gcn_input_v3}

\subsubsection{健康状态残差}

对于第$i$个状态参数，在时刻$t$的DCS实测值记为$y_i(t)$，第二章健康状态模型输出记为$\hat y_i(t)$，两者之间的状态偏差定义为
\begin{equation}
    e_i(t)=y_i(t)-\hat y_i(t).
\end{equation}
当机组处于正常状态时，GRU模型能够根据主蒸汽边界、负荷及操作量解释大部分正常状态迁移，$e_i(t)$主要反映正常预测误差和测量扰动；当设备状态发生异常后，实测状态相对于当前边界下健康参考值发生系统性偏离，多个关联参数的残差幅值、方向及动态形态将发生协同变化。因此，相比直接使用原始状态量，残差更接近“当前设备状态中健康模型无法解释的部分”。

不同压力、温度参数的原量纲及健康预测误差范围不同，为避免数值尺度影响，使用健康训练数据残差统计量进行标准化：
\begin{equation}
    r_i(t)=\frac{e_i(t)-\mu_i^{H}}{\sigma_i^{H}+\varepsilon},
\end{equation}
其中$\mu_i^{H}$和$\sigma_i^{H}$分别为健康训练集上第$i$个状态参数残差的均值和标准差，$\varepsilon$为极小常数。验证集和测试集均不得参与上述统计量计算。

\subsubsection{滑动残差窗口}

故障演化具有连续时间特征，因此以连续$L_r$个采样时刻的标准化残差构造参数--时间矩阵
\begin{equation}
R_t=
\begin{bmatrix}
 r_1(t-L_r+1)&\cdots&r_1(t)\\
 r_2(t-L_r+1)&\cdots&r_2(t)\\
 \vdots&\ddots&\vdots\\
 r_N(t-L_r+1)&\cdots&r_N(t)
\end{bmatrix}
\in\mathbb R^{N\times L_r}.
\end{equation}
其中每一行对应一个固定DCS状态节点，每一列对应一个采样时刻。本文第二章GRU健康模型的输入窗口长度与此处故障残差窗口属于两个独立超参数，前者用于建立健康动态映射，后者用于描述故障残差演化，不强制取相同数值。

在进入图卷积之前，首先利用线性映射将每个节点的时间窗口编码为$d$维初始节点特征
\begin{equation}
    H^{(0)}=\phi(R_tW_e+b_e),
\end{equation}
其中$W_e\in\mathbb R^{L_r\times d}$。随后通过归一化邻接矩阵$\bar A$进行邻域聚合，得到具有知识拓扑约束的高阶节点表示。

\subsection{面向固定测点子图的node-aware KG-GCN}
\label{subsec:node_aware_v3}

传统图分类常采用全局平均池化，将全部节点表示压缩为一个与节点排列无关的图级向量。该方式适合节点集合本身具有置换不变性的图分类任务，但本文VHP参数图中的节点与固定DCS测点一一对应，节点位置具有明确工程含义。例如，一级抽汽逆止门前后温度和1号高加入口抽汽参数虽然都属于一级抽汽支路，但其物理位置不同；VHP叶片级压力与一次冷再接口温度所代表的故障传播区域也不同。因此，如果直接计算
\begin{equation}
    \mathbf g_{mean}=\frac{1}{N}\sum_{i=1}^{N}H^{(L)}_{i,:},
\end{equation}
可能将不同局部位置出现的相似幅值异常平均为近似的图级表示。

为保留固定测点身份，本文将最后一层GCN输出按照冻结的DCS节点顺序展开：
\begin{equation}
    \mathbf g_{node}=\mathrm{vec}\left(H^{(L)}\right)
    \in\mathbb R^{Nd_L}.
\end{equation}
随后利用全连接分类头进行特征压缩并输出故障概率
\begin{equation}
    \mathbf z=\phi(W_r\mathbf g_{node}+b_r),
\end{equation}
\begin{equation}
    \hat{\mathbf p}=\mathrm{Softmax}(W_c\mathbf z+b_c).
\end{equation}
该结构仍然先利用KG邻接矩阵完成节点间信息传播，但在图级读出阶段保留各节点的空间身份，适用于本文这种“节点集合固定、测点物理位置固定”的工业参数子图。第五章将通过Mean readout与node-aware readout的对比验证这一结构选择。

分类阶段采用交叉熵损失函数
\begin{equation}
    \mathcal L_{cls}=-\frac{1}{M}\sum_{n=1}^{M}\sum_{k=1}^{C}y_{nk}\log\hat p_{nk},
\end{equation}
其中$C$为类别数，$M$为训练样本数量。

\subsection{无标签掩码预训练扩展}
\label{subsec:mask_pretrain_v3}

参考课题组既有KG-GCN方法\cite{YangHaodong2026CFB}，本文进一步设置无标签残差窗口掩码重构实验，以考察在故障标签有限条件下预训练对图编码器的影响。设无标签残差窗口为$R$，随机掩码矩阵为$M$，则掩码输入可表示为
\begin{equation}
    \tilde R=M\odot R+(1-M)\odot R_{mask}.
\end{equation}
图编码器根据$\bar A$对掩码后的节点状态进行聚合，并由解码器重构原始残差窗口
\begin{equation}
    \hat R=D_{\phi}\left(E_{\theta}(\bar A,\tilde R)\right).
\end{equation}
预训练损失仅在被遮蔽位置计算：
\begin{equation}
    \mathcal L_{mask}=\frac{\sum_{ij}(1-M_{ij})(R_{ij}-\hat R_{ij})^2}{\sum_{ij}(1-M_{ij})+\varepsilon}.
\end{equation}
需要说明的是，掩码预训练在本文中作为独立增强策略进行验证，而不预设其一定提高分类性能。其是否采用、是否冻结编码器以及最优掩码方式均由第五章验证集实验决定。这样既保留与既有KG-GCN技术路线的可比性，也避免在本文半物理故障场景下未经验证地固定训练策略。

\subsection{本章小结}

本章在第二章健康状态模型和第三章汽轮机故障知识图谱基础上，构建了面向汽轮机系统的KG-GCN故障诊断方法。首先利用健康模型实测--预测残差构造多参数动态状态特征，并从系统级知识图谱中抽取可与DCS数值映射的参数子图；随后通过归一化邻接矩阵约束GCN邻域信息传播。针对VHP代表性算例中固定DCS节点身份具有明确物理意义的特点，本文进一步采用node-aware图级读出方式，在保留知识拓扑聚合能力的同时避免全局平均池化造成局部故障位置混淆。最后设置掩码重构预训练作为少标签条件下的扩展实验。第五章将在统一半物理故障场景和多随机种子设置下，对CNN、AE、不同图拓扑以及不同读出方式进行系统比较。